\topictitle{\tr{Limite}{Limits}{Grenzwerte}}

\tr{Exemplu:}{Example:}{Beispiel:}\quad $x_n=\dfrac{2n+1}{3n+1}\ \imp\ \lim x_n=\dfrac23$

\subsection{\tr{Cazuri de \textcolor{accent}{nedeterminare}}{\textcolor{accent}{Indeterminate} forms}{\textcolor{accent}{Unbestimmte} Ausdrücke}}
\[
0\cdot(\pm\infty)\qquad\infty-\infty\qquad1^\infty\qquad0^0\qquad\infty^0\qquad\frac00\qquad\frac{\pm\infty}{\pm\infty}
\]

\subsection{1)\ \tr{Limita unui polinom}{Limit of a polynomial}{Grenzwert eines Polynoms}}
\[
\lim\bigl(a_kn^k+a_{k-1}n^{k-1}+\dots+a_1n+a_0\bigr)=\lim a_kn^k\qquad(a_k\ne0)
\]
\tr{(termenul cu $n$ la puterea cea mai mare)}{(the term with the highest power of $n$)}{(der Term mit der höchsten Potenz von $n$)}

\subsection{2)\ \tr{Limita unui cât de polinoame}{Limit of a quotient of polynomials}{Grenzwert eines Quotienten von Polynomen}}
\[
\lim\frac{a_kn^k+\dots+a_1n+a_0}{b_ln^l+\dots+b_1n+b_0}=\lim\frac{a_kn^k}{b_ln^l}
=\lim\frac{a_k}{b_l}\,n^{k-l}=
\begin{cases}0,&k<l\\[2pt]\dfrac{a_k}{b_l},&k=l\\[4pt]\dfrac{a_k}{b_l}\,(\pm\infty)^{k-l},&k>l\end{cases}
\]

\subsection{3)\ \tr{Limita unei puteri}{Limit of a power}{Grenzwert einer Potenz}}
\[
\lim a^n=\begin{cases}\infty,&a>1\\1,&a=1\\0,&a\in(-1,1)\\\nexists,&a\le-1\end{cases}
\]
\begin{remark}
\tr{Scoatem factor comun pe $a^n$, $a$ cel mai mare. Dacă avem $a^n+7^n$, discutăm cazurile $a>7$, $a<7$, $a=7$.}
    {Factor out $a^n$ with $a$ the largest base. For $a^n+7^n$ we distinguish the cases $a>7$, $a<7$, $a=7$.}
    {Man klammert $a^n$ mit der größten Basis $a$ aus. Bei $a^n+7^n$ unterscheidet man $a>7$, $a<7$, $a=7$.}
\end{remark}

\subsection{4)\ \tr{Limitele unor sume / produse}{Limits of sums / products}{Grenzwerte von Summen / Produkten}}
\tr{Cu formulele:}{With the formulas:}{Mit den Formeln:}
\[
\sum_{k=1}^{n}k=\frac{n(n+1)}{2}\qquad
\sum_{k=1}^{n}k^2=\frac{n(n+1)(2n+1)}{6}\qquad
\sum_{k=1}^{n}k^3=\frac{n^2(n+1)^2}{4}
\]

\subsection{5)\ \tr{Limitele unor șiruri cu radicali}{Limits of sequences with radicals}{Grenzwerte von Folgen mit Wurzeln}}
\begin{itemize}[label=$\blacktriangleright$]
\item $\lim\bigl(\sqrt a-b\bigr)$ \tr{(nedeterminare $\infty-\infty$)}{(form $\infty-\infty$)}{(Form $\infty-\infty$)} $\to$ \tr{raționalizăm}{rationalise}{rationalisieren}
\item $\lim\bigl(\sqrt a+b\bigr)$ $\to$ \tr{scoatem factor comun}{factor out}{ausklammern}
\end{itemize}
\begin{important}
\[
\lim\Bigl(\sqrt[3]{n^3+\cdots}-\sqrt{n^2+\cdots}\Bigr)=\lim\Bigl(\sqrt[3]{n^3+\cdots}-n\Bigr)+\lim\Bigl(n-\sqrt{n^2+\cdots}\Bigr)
\]
\end{important}

\subsection{6)\ \tr{Limite remarcabile}{Standard limits}{Wichtige Grenzwerte}}
\tr{Pentru $x_n\to0$:}{For $x_n\to0$:}{Für $x_n\to0$:}
\begin{formulae}
\item $\displaystyle\lim\frac{\sin x_n}{x_n}=1$
\item $\displaystyle\lim\frac{\arcsin x_n}{x_n}=1$
\item $\displaystyle\lim\frac{\tg x_n}{x_n}=1$
\item $\displaystyle\lim\frac{\arctg x_n}{x_n}=1$
\item $\displaystyle\lim\bigl(1+x_n\bigr)^{1/x_n}=e$\quad (\tr{forma}{form}{Form} $1^\infty$)
\item $\displaystyle\lim\frac{a^{x_n}-1}{x_n}=\ln a$
\item $\displaystyle\lim\frac{\ln(1+x_n)}{x_n}=1$
\item $\displaystyle\lim_{y_n\to1}\frac{\ln y_n}{y_n-1}=1$\quad ($\ln1=0$)
\item $\displaystyle\lim\frac{n^k}{a^n}=0$\quad ($a>1,\ k\in\N$)
\end{formulae}

\subsection{7)\ \tr{Criteriul cleștelui}{Squeeze theorem}{Einschließungssatz}}
\textit{\tr{La o sumă:}{For a sum:}{Bei einer Summe:}}
\begin{itemize}[label=$\bullet$]
\item \tr{încadrăm fiecare termen al sumei între cel mai mic și cel mai mare termen al sumei:}{bound each term of the sum between the smallest and the largest term of the sum:}{jeden Summanden zwischen dem kleinsten und dem größten Summanden einschließen:}
  \[\text{min}\le T_1\le\text{max},\quad\text{min}\le T_2\le\text{max},\ \dots,\ \text{min}\le T_n\le\text{max}\]
\item \tr{adunăm:}{add up:}{aufaddieren:}\quad $n\cdot\text{min}\le\text{sum}\le n\cdot\text{max}$
\item \tr{ambele margini tind la \dots{} $\imp$ suma tinde la \dots}{both bounds tend to \dots{} $\imp$ the sum tends to \dots}{beide Schranken streben gegen \dots{} $\imp$ die Summe strebt gegen \dots}
\end{itemize}
\begin{important}
\tr{Partea întreagă o încadrăm între:}{Bound the integer part by:}{Die Gaußklammer wird eingeschlossen durch:}
\quad $x-1<\ip{x}\le x$
\end{important}

\subsection{8)\ \tr{Cesàro--Stolz}{Cesàro--Stolz}{Cesàro--Stolz}}
\[
\lim\frac{a_n}{b_n}=\lim\frac{a_{n+1}-a_n}{b_{n+1}-b_n}
\]
\tr{dacă limita din dreapta există și $(b_n)$ este strict monotonă și nemărginită (cazul $\tfrac{\ast}{\infty}$), sau $a_n\to0$, $b_n\to0$, $(b_n)$ strict monotonă (cazul $\tfrac00$).}
    {if the right-hand limit exists and $(b_n)$ is strictly monotone and unbounded (case $\tfrac{\ast}{\infty}$), or $a_n\to0$, $b_n\to0$, $(b_n)$ strictly monotone (case $\tfrac00$).}
    {falls der rechte Grenzwert existiert und $(b_n)$ streng monoton und unbeschränkt ist (Fall $\tfrac{\ast}{\infty}$), oder $a_n\to0$, $b_n\to0$ und $(b_n)$ streng monoton (Fall $\tfrac00$).}

\subsection{9)\ \tr{Cauchy--d'Alembert (criteriul radicalului)}{Cauchy--d'Alembert (root test)}{Cauchy--d'Alembert (Wurzelkriterium)}}
\[
\lim\sqrt[n]{a_n}=\lim\frac{a_{n+1}}{a_n}\qquad(a_n>0)
\]
\tr{(dacă limita din dreapta există)}{(if the right-hand limit exists)}{(falls der rechte Grenzwert existiert)}

\subsection{10)\ \tr{Criteriul raportului}{Ratio test}{Quotientenkriterium}}
\[
\lim\frac{a_{n+1}}{a_n}=l\qquad\begin{cases}l>1&\imp\lim a_n=\infty\\ l\in[0,1)&\imp\lim a_n=0\end{cases}
\]

\subsection{11)\ L'H\^opital}
\[
\lim_{x\to x_0}\frac{f(x)}{g(x)}\ \overset{\left[\frac00\ \mt{sau}{or}{oder}\ \frac\infty\infty\right]}{=}\ \lim_{x\to x_0}\frac{f'(x)}{g'(x)}
\]
\textbf{\tr{Cazuri de nedeterminare}{Indeterminate forms}{Unbestimmte Ausdrücke}:}
\begin{itemize}[label=$\bullet$]
\item $0\cdot\infty$:\quad $f\cdot g=\dfrac{f}{1/g}\ \ \mt{sau}{or}{oder}\ \ \dfrac{g}{1/f}$,\quad \tr{apoi}{then}{dann} L'H\^opital
\item $1^\infty,\ 0^0,\ \infty^0$:\quad $f^g=e^{g\ln f}$,\quad \tr{apoi}{then}{dann} $\lim g\ln f$
\item $\infty-\infty$:\quad $f-g=f\Bigl(1-\dfrac gf\Bigr)$,\quad \tr{apoi}{then}{dann} $\lim\dfrac gf$\quad(\tr{calculăm dacă merge}{if it works}{falls es funktioniert})
\end{itemize}
\begin{important}
\[
a_n=1+\frac12+\frac13+\dots+\frac1n-\ln n\ \imp\ \lim a_n=\gamma\approx0{,}577\qquad
\bigl(\mt{constanta lui Euler}{Euler's constant}{Euler-Mascheroni-Konstante}\bigr)
\]
\end{important}

\subsection{\tr{Limite de funcții}{Limits of functions}{Grenzwerte von Funktionen}}
\begin{formulae}[1]
\item $\lim\bigl(f(x)+g(x)\bigr)=\lim f(x)+\lim g(x)$
\item $\lim\bigl(f(x)\cdot g(x)\bigr)=\lim f(x)\cdot\lim g(x)$
\item $\lim\dfrac{f(x)}{g(x)}=\dfrac{\lim f(x)}{\lim g(x)}$
\item $\lim f(x)^{g(x)}=\bigl(\lim f(x)\bigr)^{\lim g(x)}$
\item $\lim|f(x)|=\bigl|\lim f(x)\bigr|$
\end{formulae}
\begin{itemize}[label=$\blacktriangleright$]
\item \tr{Dacă un șir are limită \textbf{finită} $\imp$ șirul este \textcolor{accent}{\textbf{convergent}}.}{If a sequence has a \textbf{finite} limit $\imp$ it is \textcolor{accent}{\textbf{convergent}}.}{Hat eine Folge einen \textbf{endlichen} Grenzwert $\imp$ ist sie \textcolor{accent}{\textbf{konvergent}}.}
\item \tr{Dacă are limită \textbf{infinită} sau nu are limită $\imp$ șirul este \textcolor{accent}{\textbf{divergent}}.}{If its limit is \textbf{infinite} or does not exist $\imp$ it is \textcolor{accent}{\textbf{divergent}}.}{Ist der Grenzwert \textbf{unendlich} oder existiert nicht $\imp$ ist sie \textcolor{accent}{\textbf{divergent}}.}
\end{itemize}
